\section[Structural Criteria and Metrological Chart of \(G\)]{Criteria for selecting the gravitational coordinate: enriched holonomy, spectral
response, and metrological chart}
\label{sec:selector-gravitatorio-rev11}
\index{gravitational constant!transducer and selector}
\index{gravity!metric coordinate}
\index{holonomy!gravitational selector}

The APP--TRIT--TPK composition, the enriched state and the action branch
first construct the gravitational coordinate. The gravitational constant is
organized into three typed levels: the dimensional line that houses \(G\), the
transducer that converts a dimensionless coordinate into a section of that
line and the structural selector of that coordinate. The
Palatini--Cartan--Holst action and its constrained Hamiltonian then realize the
already generated output. The decimal chart and experimental uncertainty are
presented at the end as representation and comparison.

\subsection{Dimensional line and internal transducer}

In the dimensional lattice of Dimensional Mechanics,
\[
 \mathbf G=-\mathbf S+5\mathbf C+2\mathbf T,
 \qquad [G]=L^3M^{-1}T^{-2}.
\]
Let
\[
 c_{\rm int}\in\mathcal L_C^+,
 \qquad t_0\in\mathcal L_T^+,
 \qquad \ell_0=c_{\rm int}t_0,
 \qquad \hbar_{\rm int}\in\mathcal L_S^+.
\]

\HMTClaim{MD-R-G-TRANSDUCER}
\begin{definition}[Gravitational transducer]
\label{def:transductor-gravitatorio-rev11}
For \(\theta_G>0\), define
\begin{equation}
 \boxed{
 G_{\rm int}(\theta_G)
 =\theta_G\frac{c_{\rm int}^{5}t_0^2}{\hbar_{\rm int}}
 =\theta_G\frac{c_{\rm int}^{3}\ell_0^2}{\hbar_{\rm int}}
 \in\mathcal L_G^+.}
 \label{eq:transductor-gravitatorio-rev11}
\end{equation}
\end{definition}

The equality of the two expressions and the type of
\eqref{eq:transductor-gravitatorio-rev11} are exact. If
\[
 L_G^2:=\frac{\hbar_{\rm int}G_{\rm int}}{c_{\rm int}^3},
\]
then
\begin{equation}
 \theta_G=\left(\frac{L_G}{\ell_0}\right)^2.
 \label{eq:thetaG-razon-longitudes-rev11}
\end{equation}
The gravitational coordinate is therefore a quadratic ratio of lengths
or a normalized holonomy response. The decade
\(10^{-34}\mathcal U_S\) already belongs to \(\hbar_{\rm int}\): it fixes the absolute
action section and does not constitute a relative factor added to \(G\).

\subsection{Observable torsion and positive characters}

The observable return of \(27\) steps
\[
 \mathcal K_{27}:=F_{\rm obs}^{27}
\]
satisfies
\[
 \mathcal K_{27}^{4}=F_{\rm obs}^{108}=I.
\]
In coordinate \(\theta\in\mathbb Z/108\mathbb Z\), the order is
exactly four, since
\(\mathcal K_{27}(\theta)=\theta+27\),
\(\mathcal K_{27}^{2}(\theta)=\theta+54\neq\theta\), and
\(\mathcal K_{27}^{4}(\theta)=\theta\).

\HMTClaim{MD-R-G-POSITIVE-CHARACTER-NOGO}
\begin{theorem}[Triviality of the Observable Positive Character]
\label{thm:caracter-positivo-G-rev11}
Every multiplicative homomorphism
\[
 \chi:\langle\mathcal K_{27}\rangle
 \longrightarrow\mathbb R_{>0}^{\times}
\]
cumple \(\chi(\mathcal K_{27})=1\).
\end{theorem}

\begin{proof}
Multiplicativity implies
\(\chi(\mathcal K_{27})^4=\chi(I)=1\).
The only positive real number whose fourth power is one is \(1\).
\end{proof}

The same cycle does carry the unitary characters.
\[
 \chi_j(\mathcal K_{27})=\exp(2\pi ij/4),
 \qquad j=0,1,2,3.
\]
The observable quotient admits phase, whereas its torsion eliminates a
nontrivial positive multiplicative scale. A multiplicative continuation must therefore
act on a lift \(\widetilde{\mathcal K}_{27}\) whose class in
the abelianization of the isotropy group is nontorsional, and must naturally select
both the character and its normalization.

\subsection{Chronological memory and spectral response}
\label{sec:respuesta-espectral-G-rev11}

The directed chronological covering realizes the return of \(27\) steps as the
translation
\[
 (G_{\rm or},S)\longmapsto(G_{\rm or},S)+(1,18),
\]
of infinite semigroup order. This route preserves accumulation, but its
spectral response requires interference among histories.

Let \(\Gamma=(V,E)\) be a finite directed multigraph with positive measure,
weights \(a_\ast(e)\geq0\), and cocycle \(c(e)\in\mathbb Z^2\). The twisted
operator
\[
 (\mathcal L_{\beta,\vartheta}f)(y)
 =\sum_{e:x\to y}
 e^{-\beta a_\ast(e)}e^{i\langle\vartheta,c(e)\rangle}f(x)
\]
induce the operator positive
\(\mathcal H_{\beta,\vartheta}
=\mathcal L_{\beta,\vartheta}^{*}\mathcal L_{\beta,\vartheta}\).

\HMTClaim{MD-R-G-PHASE-INTERFERENCE}
\begin{proposition}[Interference criterion for the phase response]
\label{prop:interferencia-selector-G-rev11}
In a deterministic orbit with a unique incoming edge at each vertex,
\(\mathcal H_{\beta,\vartheta}\) is independent of \(\vartheta\). Every nonzero-phase spectral
response requires multiple coherent histories or a magnetic differential
whose holonomy survives in the enriched cycles.
\end{proposition}

\begin{proof}
On the deterministic orbit, the sole phase of each term is multiplied by its
conjugate in
\(\mathcal L_{\beta,\vartheta}^{*}\mathcal L_{\beta,\vartheta}\)
and cancels. The absence of cross terms removes every dependence on
\(\vartheta\).
\end{proof}

If a simple band \(\lambda_0(\vartheta)\) has a minimum at the origin,
its Hessian \(Q\) is positive semidefinite. With
\[
 v_{27}=(1,18),
 \qquad
 \mathcal R_{27}=v_{27}^{\mathsf T}Qv_{27},
\]
the condition \(\mathcal R_{27}\neq0\) constitutes the minimal control of
sensitivity to accumulated memory. A natural normalization
\(\mathfrak N\) would allow one to define
\[
 \Lambda_G=\ell_0\mathfrak N(Q,v_{27}),
 \qquad
 G_{\rm spec}
 =\frac{c_{\rm int}^{3}}{\hbar_{\rm int}}
  \bigl(\delta_\theta\Lambda_G\bigr)^2,
 \qquad
 \delta_\theta=\frac{A_{\rm deg}-C_{\rm deg}^{\ast}}{360}.
\]
The homogeneity of this template is exact. Its physical realization is
characterized by four positive conditions: an enriched multigraph,
a nonzero response, stability under refinement, and independence under rescaling,
together with a normalization \(\mathfrak N\) fixed before the
metrological contrast.

\subsection{Geometric Functional on Enriched Holonomy}

The second continuation defines directly a length functional.
\[
 \mathscr L:
 \operatorname{Loop}(\mathcal X_{\rm TPK}^{\rm enr})
 \longrightarrow\mathcal L_L^+\cup\{0\},
\]
invariant under cyclic change of basepoint, conjugation, admissible
relabeling, and subdivision. For an enriched lift, write
\[
 \Lambda_{27}:=\mathscr L(\widetilde{\mathcal K}_{27}),
 \qquad
 \rho_{27}:=\frac{\Lambda_{27}}{c_{\rm int}t_0},
 \qquad
 \theta_G^{\rm hol}:=(\delta_\theta\rho_{27})^2.
\]
Transduction produces the homogeneous template
\begin{equation}
 \boxed{
 G_{\rm hol}
 =\frac{c_{\rm int}^{3}}{\hbar_{\rm int}}
  \bigl(\delta_\theta\Lambda_{27}\bigr)^2.}
 \label{eq:plantilla-holonomia-G-rev11}
\end{equation}
The concatenation law of this functional on the torsional cycle forms
part of its definition and its proof. Physical selection is completed
by constructing \(\widetilde{\mathcal K}_{27}\) and \(\mathscr L\), fixing a
unique normalization, and verifying the stability of the output. The value
\(G_{\rm CODATA\,2022}\) enters exclusively afterward, in the metrological
comparison.

\begin{proposition}[Separation of the two gravitational selectors]
\label{prop:separacion-selectores-G}
The circular selector
\(\theta_{108}^{\rm circ}=(54/\pi)^2\), constructed by action--Compton--gravity
closure, and the holonomic selector
\(\theta_G^{\rm hol}=(\delta_\theta\rho_{27})^2\) are two distinct typed realizations
of the transducer \(G_{\rm int}\).  They are not identified with
one another, and the decimal chart \(\mathscr R_G\) does not constitute an evaluation of
\(G_{\rm hol}\).
\end{proposition}
\begin{proof}
The first selector depends solely on the circular closure of order 108. The second additionally contains the spectral length \(\Lambda_{27}=\mathscr L(\widetilde{\mathcal K}_{27})\). An equality between them would require computing \(\Lambda_{27}\) and proving the agreement of normalizations; no operation of the decimal reader performs that computation. Their domains and input data therefore distinguish them before any numerical comparison.
\end{proof}

\subsection{Structural selection criterion}

\HMTClaim{MD-R-G-SELECTOR-CRITERION}
\begin{criterion}[Structural selector of the gravitational constant]
\label{crit:selector-G-rev11}
An HMT selector of \(G\) is identified when it jointly satisfies:
\begin{enumerate}
 \item enriched domain, composition, and orientation declared;
 \item descent to a nontorsional Abelian class, or a geometric functional
 invariant under basepoint, conjugation, relabelling, and subdivision;
 \item nontrivial positive response stable under refinement;
 \item unique normalization independent of the operator's overall scale;
 \item internal selection of the coordinate and its decimal position;
 \item freezing of the code and entries before revealing the
 comparison value;
 \item final realization as a section of \(\mathcal L_G\) in a common
 dimensional chart.
\end{enumerate}
\end{criterion}

The dimensionless ratios
\[
 \frac{\mathcal G_{\rm tor}}{\mathcal I_{\rm tor}}=1,
 \qquad
 \frac{\mathcal G_{\rm av}}{\mathcal I_{\rm av}}
 =\frac{\pi_{\rm HMT}}{\varphi^2}
\]
describe distinct gravitational--inertial readers and remain
compatible with a single dimensional section \(G\). The construction of that
section also uses
\((\hbar_{\rm int},c_{\rm int},\ell_0)\) and the selector \(\theta_G\).

\subsection{Evaluation of the decimal gravitational chart and subsequent metrological contrast}

The decimal representation preserved in the corpus is
\begin{equation}
 \mathscr R_G([n,k,d]_3;t,e)
 :=10^{-n}\left(k+\frac{t}{10^3}+\frac{e}{10^d}\right).
 \label{eq:lector-decimal-G-rev11}
\end{equation}
For the chart
\([n,k,d]_3=[11,6,5]_3\), \(t=674\) and \(e=30\), the evaluation racional is
\begin{equation}
 \boxed{
 \mathscr R_G
 =10^{-11}\left(6+\frac{674}{1000}+\frac{30}{10^5}\right)
 =6.67430\times10^{-11}.}
 \label{eq:evaluacion-decimal-G-rev11}
\end{equation}
The triad, the two blocks, and the decimal position constitute the declared
data of this chart. Their evaluation is exact; their structural selection
is governed by \cref{crit:selector-G-rev11}.

After choosing the SI basis of \(\mathcal L_G\), the current contrast is
\begin{equation}
 \boxed{
 G_{\rm CODATA\,2022}=6.67430(15)\times10^{-11}
 \ \mathrm{m^3\,kg^{-1}\,s^{-2}}.}
 \label{eq:G-CODATA-2022-rev11}
\end{equation}

\begin{center}
\begin{tabularx}{0.94\textwidth}{@{}L{0.27\textwidth}Y@{}}
\toprule
Field & Metrological Statement\\
\midrule
Central value &
\(6.67430\times10^{-11}\ \mathrm{m^3\,kg^{-1}\,s^{-2}}\)\\
Standard uncertainty &
\(u(G)=0.00015\times10^{-11}
=1.5\times10^{-15}\ \mathrm{m^3\,kg^{-1}\,s^{-2}}\)\\
Incertidumbre relativa & \(u_r(G)\simeq2.2\times10^{-5}\)\\
Source and temporality &
2022 CODATA recommended values, published in 2024
\cite{CODATA2022}\\
Proof function &
Subsequent external contrast; neither the unit nor the uncertainty enters into
the rational evaluation of \eqref{eq:evaluacion-decimal-G-rev11}.\\
\bottomrule
\end{tabularx}
\end{center}

\paragraph{Reproducibility of the gravitational reader.}
An independent enumeration jointly verifies the dimensional type, the
decimal reader, and the monodromy identities on the declared finite
domain.

\begin{proposition}[Causal chain of the gravitational constant]
The covariant action, the restricted Hamiltonian, and the
Einstein--Schrödinger equation govern the gravitational dynamics. The dimensional type
and the transducer \eqref{eq:transductor-gravitatorio-rev11} are exact; the
evaluation \eqref{eq:evaluacion-decimal-G-rev11} is exact for its declared
chart. The value, unit, and uncertainty of
\eqref{eq:G-CODATA-2022-rev11} constitute an external comparison with CODATA 2022. The
autonomous selection of \(\theta_G\) is typed by the enriched
lift, the interfering response, and the holonomy functional subject
to the \cref{crit:selector-G-rev11}.
\end{proposition}
